% !TeX root = ../../Book3.tex
% 纲要标题：### 16.3 深度引理与合冲模
% 纲要位置：工作记录/计划纲要.md，第 2501 行。
% 纲要细目（写作范围）：
% * 短正合列中的深度不等式与零模的无穷大约定
% * 合冲模的深度下界及其不能充当上界的例子
% * 用消解逐步比较深度

\section{深度引理与合冲模}
\label{b3:sec:1603}

短正合列把三个模联系起来，深度却通常只满足不等式。首次非零的 Ext 及其连接同态可以说明何时达到等号，也能追踪商去一个非零因子后深度的变化。


% 纲要内容（仅作写作计划，尚非教材正文）：
%
% * 短正合列中的深度不等式与零模的无穷大约定
% * 合冲模的深度下界及其不能充当上界的例子
% * 用消解逐步比较深度
%

\subsection{短正合列中的深度不等式}
\label{b3:subsec:1603:01}

\begin{theorem}[深度引理]\label{b3:thm:depth-lemma}
设 \((R,\mathfrak m)\) 为交换 Noether 局部环，\(k=R/\mathfrak m\) 为其剩余域，
\[
0\longrightarrow A\longrightarrow B\longrightarrow C\longrightarrow0
\]
为有限生成的 \(R\) 模的短正合列。约定零模深度为 \(+\infty\)，并采用
\[
(+\infty)+1=(+\infty)-1=+\infty,
\qquad \min\{+\infty,t\}=t
\quad(t\in\mathbb Z\cup\{+\infty\}).
\]
则
\begin{align*}
\operatorname{depth}B&\ge\min\{\operatorname{depth}A,\operatorname{depth}C\},\\
\operatorname{depth}A&\ge\min\{\operatorname{depth}B,\operatorname{depth}C+1\},\\
\operatorname{depth}C&\ge\min\{\operatorname{depth}A-1,\operatorname{depth}B\}.
\end{align*}
此外，若 \(\operatorname{depth}B>\operatorname{depth}C\) 且 \(C\ne0\)，则 \(\operatorname{depth}A=\operatorname{depth}C+1\)；若 \(\operatorname{depth}A<\operatorname{depth}B\) 且 \(A\ne0\)，则 \(\operatorname{depth}C=\operatorname{depth}A-1\)。
\end{theorem}
\begin{proof}
令 \(E^i(X)=\operatorname{Ext}_R^i(k,X)\)，并对 \(i<0\) 置 \(E^i(X)=0\)。短正合列给出
\[
\cdots\longrightarrow E^{i-1}(C)\longrightarrow E^i(A)
\longrightarrow E^i(B)\longrightarrow E^i(C)
\longrightarrow E^{i+1}(A)\longrightarrow\cdots.
\]
由\cref{b3:thm:depth-ext}，深度就是首次非零次数。三条不等式分别读这条长列中以 \(E^i(B)\)、\(E^i(A)\)、\(E^i(C)\) 为中项的三个片段。若 \(i\) 小于 \(A,C\) 的两个深度，\(E^i(A)=E^i(C)=0\)，故 \(E^i(B)=0\)，得到第一式。若 \(i\) 小于 \(\operatorname{depth}B\) 与 \(\operatorname{depth}C+1\)，则 \(E^{i-1}(C)=E^i(B)=0\)，得到第二式；同样由 \(E^i(B)=E^{i+1}(A)=0\) 得第三式。

若 \(c=\operatorname{depth}C<\operatorname{depth}B\)，则非零模 \(E^c(C)\) 单射到 \(E^{c+1}(A)\)。第二式又给出 \(\operatorname{depth}A\ge c+1\)，故恰等于 \(c+1\)。若 \(a=\operatorname{depth}A<\operatorname{depth}B\)，则长列使 \(E^{a-1}(C)\) 满射到非零的 \(E^a(A)\)；这也表明 \(a\ne0\)。第三式给出反向界，故深度恰为 \(a-1\)。这两个等式都来自连接映射：中项尚未出现非零 Ext 时，一端的首次非零类迫使另一端在相邻次数非零。零模的情形可从同构直接核对。
\end{proof}

三个下界都不能普遍改成等式。例如分裂列 \(0\rightarrow A\rightarrow A\oplus C\rightarrow C\rightarrow0\) 中，\(B\) 的深度确为两端较小者；但设 \(k\) 为域，对 \(R=k[X,Y]_{(X,Y)}\) 的列 \(0\rightarrow R\xrightarrow{X}R\rightarrow R/(X)\rightarrow0\)，两端深度分别为二和一，中项深度为二。

\subsection{合冲模的深度}
\label{b3:subsec:1603:02}

\begin{definition}\label{b3:def:syzygy-local}
设 \(R\) 为交换 Noether 环，\(M\) 为有限生成的 \(R\) 模，给定有限秩自由模 \(F\) 到 \(M\) 的满射。其核 \(\Omega M=\ker(F\rightarrow M)\) 称为这次自由展示所得的\term[hechongmo]{合冲模}[syzygy module]。重复对核选择自由满射，得到更高合冲模 \(\Omega^jM\)，其中 \(\Omega^0M=M\)。
\end{definition}

合冲模依赖所选展示。例如把 \(F\twoheadrightarrow M\) 换成 \(F\oplus R\twoheadrightarrow M\)，让新增分量映到零，核就从 \(\Omega M\) 变成 \(\Omega M\oplus R\)。当前记号因此总是指选定展示的核；本节的深度估计对每个选择都成立。下一章将用局部极小生成组固定更精细的消解。

\begin{corollary}\label{b3:cor:syzygy-depth}
设 \((R,\mathfrak m)\) 为交换 Noether 局部环，\(M\) 为有限生成的 \(R\) 模，\(j\ge0\) 为整数。在 \(M\) 的任一各项均为有限秩自由模的消解中，第 \(j\) 个合冲模满足
\[
\operatorname{depth}_R\Omega^jM
\ge\min\{\operatorname{depth}R,\operatorname{depth}_RM+j\}.
\]
若 \(M\ne0\)、\(\operatorname{depth}_RM<\operatorname{depth}R\)，则第一合冲模满足 \(\operatorname{depth}\Omega M=\operatorname{depth}M+1\)。
\end{corollary}
\begin{proof}
非零有限自由模 \(F=R^n\) 与 \(R\) 有相同深度，因为正则性在每个直和分量逐项检验。对 \(0\rightarrow\Omega M\rightarrow F\rightarrow M\rightarrow0\) 应用深度引理第二式，得到 \(\operatorname{depth}\Omega M\ge\min\{\operatorname{depth}R,\operatorname{depth}M+1\}\)；每换成一个更高的核，都可以重复这一步，逐次归纳即得下界。若两深度严格不等，深度引理的第一条附加等式给出精确值。遇到零合冲模时深度为无穷大，下界仍成立；这个约定不表示零模存在无限长的非退化正则序列。
\end{proof}

这个估计只是下界，\(\operatorname{depth}R\) 并不是合冲模实际深度的上界。设 \(k\) 为域。在 \(R=k[X,Y]_{(X,Y)}\) 上，剩余域 \(k\) 深度为零，其第一合冲模 \(\mathfrak m=(X,Y)\) 深度为一。第二合冲模由 \((-Y,X)\) 生成，与 \(R\) 同构，深度为二。这把前章 Koszul 消解的三个次数与深度的逐步增加对应起来。

\subsection{沿消解逐步比较深度}
\label{b3:subsec:1603:03}

当 \(\operatorname{depth}M+j\ge\operatorname{depth}R\) 时，上述下界不再随 \(j\) 增长；下面的模却会有高于环本身的深度。

\ProseParagraphBreak

例如，设 \(k\) 为域，取
\[
R=k[x,y]_{(x,y)}/(x^2,xy),\qquad M=R/(y).
\]
以下仍用 \(x,y\) 表示它们在 \(R\) 中的像。非零元 \(x\) 被极大理想 \((x,y)\) 零化，故 \(\operatorname{depth}R=0\)。自然满射 \(R\to M\) 的核为 \((y)\)。由于 \(xy=0\)，而乘 \(y\) 在 \(R/(x)\cong k[y]_{(y)}\) 上单射，\(\operatorname{Ann}_R(y)=(x)\)，所以
\[
\Omega M=(y)\cong R/(x)\cong k[y]_{(y)}.
\]
在该模上乘 \(y\) 单射，其支集维数为一，因而 \(\operatorname{depth}_R\Omega M=1>\operatorname{depth}R\)。消解长度与深度之间的精确关系还需要另外证明。

以最左端自由模的深度下界 \(\operatorname{depth}R\) 为起点，沿消解向右逐次使用深度引理，所得深度下界每步至多降低一。因此一个长度 \(n\) 的消解给出 \(\operatorname{depth}R-n\) 这个下界，也就对消解长度施加了必要限制。

\begin{proposition}\label{b3:prop:finite-resolution-depth-bound}
设 \((R,\mathfrak m)\) 为交换 Noether 局部环，\(M\ne0\) 为有限生成的 \(R\) 模，\(n\ge0\) 为整数。若有长度 \(n\)、各项均为有限秩自由模的消解
\[
0\longrightarrow F_n\longrightarrow\cdots\longrightarrow F_0
\longrightarrow M\longrightarrow0,
\]
则 \(\operatorname{depth}_RM\ge\operatorname{depth}R-n\)。
\end{proposition}
\begin{proof}
若 \(n=0\)，\(M\) 自由，结论成立。一般令 \(L=\ker(F_0\rightarrow M)\)。若 \(L=0\)，结论仍由自由性得出；否则尾部是 \(L\) 的长度 \(n-1\) 自由消解，归纳给出 \(\operatorname{depth}L\ge\operatorname{depth}R-(n-1)\)。深度引理第三式于是给出
\[
\operatorname{depth}M\ge
\min\{\operatorname{depth}L-1,\operatorname{depth}R\}
\ge\operatorname{depth}R-n.
\]
\end{proof}

例如，设 \(k\) 为域，取 \(R=k[X,Y,Z]_{(X,Y,Z)}\)、\(M=R/(X,Y)\)。Koszul 消解长为二，上式给出深度至少一；余下的 \(Z\) 确为正则元素，而商的维数为一，故深度恰为一。一般地，上式也写成 \(n\ge\operatorname{depth}R-\operatorname{depth}M\)，但这里的 \(n\) 是给定消解的长度，未必最短。对有限投射维数模，下一章将用极小消解证明 Auslander--Buchsbaum 公式，把这个必要界改进为投射维数的精确等式。
